\section{Interpretation}\label{sec:mechanisms}

\subsection{Why equities and not Treasuries: tick size as a limit to arbitrage}

The surprise explains the first-minute response of \ZN\ as well as that of the equity indices, so
the absence of a tradeable drift in \ZN\ is not a failure of the signal. Two features of the
Treasury market explain it. First, a CPI or payrolls surprise maps into the expected policy path,
and hence into Treasury yields, almost mechanically, and the Treasury futures market is where the
fastest participants price that information. By the time a position can be opened, the move is
largely complete: the drift per standard deviation of the CPI signal is 0.8 bps in \ZN, against
about 6 bps in \ES\ and 8.5 bps in \NQ. For equities, investors must also judge the effect of the same number on
growth, margins and risk premia, and the drift is largest in \NQ, whose value depends most on
long-dated cash flows; \YM, with more industrial and financial stocks, reacts about 8\% less than
\ES\ but drifts by the same amount after the first minute. Second, even where a gross edge exists, one \ZN\ tick is large relative to
the post-release move: the naive CPI rule earns 2.7 bps per trade before costs and pays 3.0 bps to
trade, and the DTW ladder's gross Sharpe ratio of 0.34 becomes $-0.97$ at a quarter tick per side.
The minimum price increment thus acts as a limit to arbitrage that leaves a small predictable
component in prices unexploited.

\subsection{Why price-path analogues fail on direction}

Analogue forecasts assume that the mapping from recent price shapes to subsequent returns is
stable. Around macro releases it was not. The analogues' direction forecast had a negative rank
correlation with subsequent returns in 2013--2017 and a positive one in 2018--2020, and the
single-window analogue book that earned money in 2018--2020 lost money in \ES\ and \NQ\ in
2021--2026; on \YM, added later, it earned little in either period. Across 567 constructions of the analogue forecast (waiting times, path
lengths and neighbour weights) the training and validation correlations are almost unrelated
(correlation 0.07 across settings). The change coincides with the shift of the macro regime from
low inflation to the 2021--2022 inflation shock, when the same price shape could precede
continuation or reversal depending on how the market read the Federal Reserve's response. What the
analogues do capture is uncertainty: when similar past paths were followed by very different
outcomes, the next move tends to be large. This is why their dispersion is a stable forecast of
the size of the move and a useful input for position sizing, while their direction is not.

\subsection{What the entry assumption reveals about the speed of price discovery}

The classifier result in Table~\ref{tab:equalcosts} quantifies how much of the forecastable
component of returns is priced within the first minute. Moving the entry by one minute removes
about 80\% of the gross profit of a model that combines surprise and price information. Together
with the event-study evidence that roughly a quarter of the CPI response remains after the first
minute in equity futures, this places most of the price discovery within the first minute and a
small, economically meaningful residual in the following half hour. The residual is predictable
from public information, persists for about 30 minutes---a second 30-minute window after the release
earns negative returns in every market---and is concentrated in periods when macro data drive
expectations of monetary policy.
